\chapter{Lezione 11 --- Progetto del controllo digitale via loop shaping continuo, filtro anti-alias, sample-and-hold e discretizzazione del regolatore}

\section{Obiettivo della lezione}

In questa lezione si imposta una procedura pratica di progetto per un sistema di controllo digitale.

L'idea di fondo è la seguente:
\begin{itemize}
    \item l'impianto fisico resta un sistema continuo;
    \item il regolatore viene implementato in digitale;
    \item per il progetto, per\`o, si sfruttano gli strumenti pi\`u ricchi del controllo continuo;
    \item solo alla fine il regolatore continuo viene discretizzato.
\end{itemize}

Questa scelta \`e molto naturale per sistemi SISO tempo-invarianti, perch\'e il progetto nel continuo,
specialmente tramite diagrammi di Bode e loop shaping, \`e pi\`u intuitivo e pi\`u ricco
dal punto di vista degli strumenti disponibili.

\section{Schema del sistema di controllo digitale}

Lo schema fisico considerato \`e il seguente:
\[
y^\star \longrightarrow \text{sommatore} \longrightarrow \text{campionatore} \longrightarrow R(z)
\longrightarrow \text{D/A} \longrightarrow \text{amplificatore + attuatore}
\longrightarrow G(s) \longrightarrow y.
\]

Prima del campionatore bisogna inoltre prevedere un filtro anti-alias analogico.

Quindi, dal punto di vista fisico e implementativo, i blocchi aggiuntivi rispetto al solo impianto sono:
\begin{itemize}
    \item \textbf{\(Q_1(s)\)}: filtro anti-alias;
    \item \textbf{\(Q_2(s)\)}: convertitore D/A, modellato come tenuta di ordine zero;
    \item \textbf{\(Q_3(s)\)}: amplificatore + attuatore.
\end{itemize}

A seconda della convenzione usata, nel modello globale pu\`o essere inclusa anche la dinamica del sensore.

\section{Modello continuo equivalente usato per il progetto}

Per il progetto si inglobano tutte le dinamiche continue nel blocco
\[
\tilde G(s)=Q_1(s)\,Q_2(s)\,Q_3(s)\,G(s).
\]

Si progetta quindi un regolatore continuo \(C(s)\) come se il sistema fosse interamente continuo,
in modo tale che la funzione di trasferimento in ciclo chiuso
\[
H(s)=\frac{C(s)\tilde G(s)}{1+C(s)\tilde G(s)}
\]
soddisfi le specifiche assegnate.

Una volta determinato \(C(s)\), lo si discretizza per ottenere \(R(z)\),
ad esempio con la trasformata di Tustin, con o senza pre-warping.

\section{Perch\'e si progetta \(C(s)\) e non direttamente \(R(z)\)}

Il motivo \`e pratico:
nel continuo si hanno a disposizione strumenti di progetto molto efficaci:
\begin{itemize}
    \item specifiche statiche;
    \item specifiche in transitorio;
    \item margine di fase;
    \item banda passante;
    \item reiezione dei disturbi;
    \item loop shaping sulla funzione d'anello.
\end{itemize}

Per questo si sceglie di:
\begin{enumerate}
    \item progettare \(C(s)\) nel continuo;
    \item poi trasformarlo in \(R(z)\).
\end{enumerate}

\section{Specifiche del problema}

L'impostazione fatta a lezione considera specifiche tipiche del controllo:
\begin{itemize}
    \item prestazioni statiche;
    \item prestazioni di transitorio;
    \item reiezione dei disturbi e del rumore di misura.
\end{itemize}

Nell'esempio svolto, il sistema continuo da controllare \`e
\[
G(s)=\frac{10}{s+3}.
\]

Le specifiche richieste per il sistema in ciclo chiuso sono del tipo:
\begin{itemize}
    \item errore nullo a regime in risposta al gradino;
    \item tempo di assestamento al \(95\%\) inferiore a \(2\) s;
    \item assenza di sovraelongazione nella risposta al gradino;
    \item reiezione del rumore di misura di almeno \(40\) dB per frequenze superiori a \(60\) rad/s.
\end{itemize}

\section{Forma della funzione d'anello}

Nel progetto per loop shaping si lavora sulla funzione di anello aperto
\[
L(s)=C(s)G(s)
\]
oppure, pi\`u correttamente, sul sistema globale continuo
\[
L(s)=C(s)\tilde G(s).
\]

La funzione in ciclo chiuso \`e
\[
H(s)=\frac{L(s)}{1+L(s)}.
\]

\section{Specifica di errore a regime nullo}

Per avere errore nullo al gradino occorre che la funzione d'anello abbia almeno un polo nell'origine.

Quindi il regolatore deve contenere un integratore:
\[
C(s)\sim \frac{1}{s}.
\]

Questa \`e la prima scelta strutturale da fare nel progetto.

\section{Specifica di assenza di sovraelongazione}

Per evitare sovraelongazione, si vuole un sistema in ciclo chiuso con comportamento fortemente smorzato,
idealmente con polo dominante reale oppure con coppia dominante molto smorzata.

Negli appunti viene richiamata la regola empirica gi\`a vista nelle lezioni precedenti:
\[
\zeta \approx \frac{M_\varphi}{100},
\]
con \(M_\varphi\) espresso in gradi.

Per avere una risposta senza overshoot, si richiede quindi un margine di fase elevato, ad esempio
\[
M_\varphi > 75^\circ.
\]

\section{Specifica sul tempo di assestamento}

Se la risposta chiusa pu\`o essere approssimata con quella di un primo ordine
\[
H(s)=\frac{1}{1+s\theta},
\]
allora la risposta al gradino \`e
\[
y(t)=1-e^{-t/\theta}.
\]

Imponendo il tempo di assestamento al \(95\%\),
\[
1-e^{-t_a/\theta}=0.95,
\]
si ricava
\[
e^{-t_a/\theta}=0.05,
\qquad
-\frac{t_a}{\theta}=\ln(0.05),
\qquad
\theta=\frac{t_a}{-\ln(0.05)}.
\]

Con la specifica
\[
t_{a,95\%}<2\ \text{s},
\]
segue
\[
\theta<\frac{2}{-\ln(0.05)} \approx 0.668.
\]

Quindi
\[
\frac{1}{\theta} \gtrsim 1.5\ \text{rad/s}.
\]

Negli appunti, la banda passante desiderata del sistema chiuso viene quindi scelta attorno a
\[
\omega_{B,-3}\approx 2\ \text{rad/s}.
\]

\section{Specifica di reiezione del rumore}

La richiesta di attenuare il rumore di misura di almeno \(40\) dB per
\[
\omega>60\ \text{rad/s}
\]
si traduce in una barriera sul modulo della funzione d'anello:
alle alte frequenze \(L(s)\) deve essere sufficientemente piccolo,
cos\`i che il rumore venga attenuato.

Questo vincolo impone che, oltre una certa pulsazione, il modulo di \(L(s)\)
scenda rapidamente.

\section{Regolatori di tentativo e loop shaping}

Negli appunti il regolatore viene costruito per tentativi, cio\`e plasmando il Bode di \(L(s)\).

\subsection{Primo tentativo}

Una prima scelta naturale \`e
\[
C_1(s)=\frac{k}{s},
\]
in modo da:
\begin{itemize}
    \item introdurre il polo nell'origine;
    \item avere elevato guadagno a bassa frequenza;
    \item iniziare con una pendenza di \(-20\) dB/dec nel diagramma del modulo.
\end{itemize}

\subsection{Correzione con una rete anticipatrice/ritardatrice}

Per soddisfare contemporaneamente:
\begin{itemize}
    \item crossover desiderato;
    \item margine di fase elevato;
    \item attenuazione ad alta frequenza;
\end{itemize}
si modifica il regolatore introducendo uno zero e un polo.

Dalle formule annotate in foto emerge la forma finale del regolatore continuo:
\[
C(s)=\frac{3}{10}\,\frac{1}{s}\,
\frac{1+\dfrac{s}{0.7}}{1+\dfrac{s}{1.5}}.
\]

Questa struttura:
\begin{itemize}
    \item mantiene l'integratore;
    \item aggiunge un aumento di fase nell'intorno della banda d'interesse;
    \item consente di posizionare il taglio a \(0\) dB circa in
    \[
    \omega_c \approx 2\ \text{rad/s};
    \]
    \item lascia poi al polo dell'impianto e agli altri poli il compito di far scendere il modulo alle alte frequenze.
\end{itemize}

\section{Interpretazione del Bode del progetto}

La logica del progetto riportata in lavagna \`e la seguente:
\begin{itemize}
    \item a bassissima frequenza serve l'integratore;
    \item vicino alla zona del crossover si introducono zeri per aumentare la fase;
    \item si lascia poi che i poli del sistema riportino il modulo verso il basso;
    \item ad alta frequenza si aggiungono poli per rispettare la barriera di attenuazione del rumore.
\end{itemize}

Questo \`e il tipico schema di loop shaping:
\[
\text{guadagno alto in bassa frequenza, fase sufficiente al crossover, modulo piccolo in alta frequenza.}
\]

\section{Scelta del periodo di campionamento}

Dopo aver fissato le prestazioni in ciclo chiuso, si passa alla scelta del campionamento.

Negli appunti si usa la regola empirica:
\[
\omega_c^\text{camp} \approx 10\,\omega_{B,-3},
\]
dove \(\omega_c^\text{camp}\) \`e la pulsazione di campionamento.

Poich\'e nel progetto si \`e scelto
\[
\omega_{B,-3}\approx 2\ \text{rad/s},
\]
una prima scelta naturale sarebbe
\[
\omega_c^\text{camp}=20\ \text{rad/s}.
\]

La pulsazione di Nyquist corrispondente sarebbe
\[
\omega_N=\frac{\omega_c^\text{camp}}{2}=10\ \text{rad/s}.
\]

\section{Scelta del filtro anti-alias}

Il filtro anti-alias analogico deve avere una banda passante non superiore alla Nyquist,
per evitare aliasing del rumore e delle componenti veloci.

Quindi, con la prima scelta
\[
\omega_N=10\ \text{rad/s},
\]
bisognerebbe scegliere un anti-alias con frequenza di taglio non superiore a tale valore.

Negli appunti si osserva per\`o che questa scelta pu\`o essere troppo penalizzante,
e si preferisce aumentare la frequenza di campionamento.

Un esempio riportato \`e scegliere
\[
\omega_c^\text{camp}=60\ \text{rad/s},
\qquad
\omega_N=30\ \text{rad/s},
\]
cos\`i da poter usare un filtro anti-alias con pulsazione caratteristica attorno a \(30\) rad/s.

Una possibile scelta di anti-alias richiamata a lezione \`e un passa-basso del terzo ordine del tipo
\[
Q_1(s)=\left(\frac{30}{s+30}\right)^3.
\]

\subsection*{Periodo di campionamento}

Con
\[
\omega_c^\text{camp}=60\ \text{rad/s},
\]
si ottiene
\[
T=\frac{2\pi}{\omega_c^\text{camp}}=\frac{2\pi}{60}\approx 0.1\ \text{s}.
\]

\section{Effetto della tenuta di ordine zero sul margine di fase}

Una volta fissato \(T\), bisogna tenere conto del convertitore D/A con tenuta di ordine zero.

Alle frequenze basse rispetto alla Nyquist, la tenuta si comporta come
\[
F(s)\approx T e^{-sT/2},
\]
quindi introduce un ritardo puro approssimato pari a \(T/2\).

Questo significa che alla pulsazione di crossover \(\omega\) si ha un contributo di fase circa
\[
-\omega \frac{T}{2}.
\]

Nel caso dell'esempio, con
\[
\omega=2\ \text{rad/s},
\qquad
T\approx 0.1\ \text{s},
\]
si ottiene
\[
-\omega\frac{T}{2}\approx -0.1\ \text{rad}\approx -6^\circ.
\]

Negli appunti questa perdita di fase viene ritenuta non trascurabile:
se il margine di fase \`e gi\`a al limite, \(6^\circ\) possono compromettere la specifica.

\section{Campionare pi\`u velocemente}

Per ridurre la perdita di fase dovuta alla tenuta, conviene campionare pi\`u rapidamente,
cio\`e diminuire \(T\).

L'idea annotata a lezione \`e quindi:
\begin{itemize}
    \item mantenere l'anti-alias fissato sulla banda scelta;
    \item aumentare ulteriormente la frequenza di campionamento;
    \item in questo modo il ritardo equivalente \(T/2\) diventa pi\`u piccolo.
\end{itemize}

Cos\`i la perdita di fase della tenuta si riduce e il margine di fase resta compatibile con la specifica.

\section{Da \(C(s)\) a \(R(z)\)}

Una volta definito il regolatore continuo \(C(s)\), si passa a costruire il regolatore digitale \(R(z)\).

La trasformazione scelta a lezione \`e la trasformata di Tustin:
\[
s=\frac{2}{T}\frac{z-1}{z+1}.
\]

\subsection{Pre-warping: serve o no?}

Negli appunti compare la domanda:
\begin{quote}
\emph{Da \(C(s)\) a \(R(z)\) faccio la trasformata di Tustin. Ma faccio o non faccio il pre-warping?}
\end{quote}

La risposta data nel caso in esame \`e:
\begin{itemize}
    \item s\`i, conviene farlo;
    \item perch\'e il progetto del margine di fase \`e gi\`a piuttosto tirato;
    \item senza pre-warping, la deformazione in frequenza della Tustin potrebbe far perdere il rispetto della specifica.
\end{itemize}

Quindi si usa la Tustin con pre-warping sulla frequenza di maggiore interesse,
tipicamente la frequenza di crossover o la frequenza di taglio desiderata.

\section{Riserva di fase nel progetto continuo}

Un messaggio importante della lezione \`e che,
nel progetto continuo del regolatore,
non conviene soddisfare la specifica sul margine di fase in modo troppo poco conservativo.

Infatti, una parte del margine viene poi ``mangiata'' da:
\begin{itemize}
    \item filtro anti-alias;
    \item tenuta di ordine zero;
    \item eventuali altre dinamiche implementative.
\end{itemize}

Per questo \`e buona norma:
\begin{itemize}
    \item o campionare abbastanza velocemente da rendere trascurabili questi effetti;
    \item oppure progettare \(C(s)\) con una riserva di fase aggiuntiva.
\end{itemize}

\section{Procedura complessiva di progetto}

La procedura emersa in questa lezione pu\`o essere riassunta cos\`i:

\begin{enumerate}
    \item Si modellano tutte le dinamiche continue presenti:
    \[
    \tilde G(s)=Q_1(s)Q_2(s)Q_3(s)G(s).
    \]

    \item Si fissano le specifiche:
    \begin{itemize}
        \item errore a regime;
        \item tempo di assestamento;
        \item overshoot;
        \item reiezione del rumore.
    \end{itemize}

    \item Si progetta un regolatore continuo \(C(s)\) mediante loop shaping,
    imponendo che
    \[
    H(s)=\frac{C(s)\tilde G(s)}{1+C(s)\tilde G(s)}
    \]
    soddisfi le specifiche.

    \item Si sceglie la banda passante desiderata del sistema chiuso.

    \item Si sceglie una frequenza di campionamento sufficientemente maggiore della banda chiusa.

    \item Si sceglie un filtro anti-alias compatibile con la Nyquist.

    \item Si verifica la perdita di fase introdotta dalla tenuta di ordine zero.

    \item Si discretizza il regolatore con Tustin, eventualmente con pre-warping.

    \item Si verifica infine il comportamento del sistema discreto tramite simulazione,
    ad esempio con risposta al gradino, rampa e analisi del margine.
\end{enumerate}

\section{Sintesi della lezione}

I punti chiave della lezione sono:
\begin{enumerate}
    \item nel progetto digitale conviene inglobare le dinamiche implementative nel modello continuo
    \[
    \tilde G(s)=Q_1(s)Q_2(s)Q_3(s)G(s);
    \]
    \item il regolatore si progetta inizialmente come \(C(s)\), usando gli strumenti del controllo continuo;
    \item nell'esempio con
    \[
    G(s)=\frac{10}{s+3},
    \]
    si richiedono errore nullo al gradino, transitorio rapido, assenza di overshoot e attenuazione del rumore;
    \item per avere errore nullo al gradino serve un polo nell'origine nel regolatore;
    \item per evitare overshoot si vuole un margine di fase elevato, dell'ordine di \(75^\circ\);
    \item un regolatore continuo compatibile con gli appunti \`e
    \[
    C(s)=\frac{3}{10}\,\frac{1}{s}\,
    \frac{1+\dfrac{s}{0.7}}{1+\dfrac{s}{1.5}};
    \]
    \item la banda passante chiusa viene scelta circa a
    \[
    \omega_{B,-3}\approx 2\ \text{rad/s};
    \]
    \item il periodo di campionamento va scelto a partire da questa banda;
    \item il filtro anti-alias e la tenuta di ordine zero introducono effetti dinamici non trascurabili;
    \item la tenuta introduce una perdita di fase circa pari a
    \[
    -\omega T/2;
    \]
    \item se questa perdita \`e rilevante, bisogna campionare pi\`u velocemente oppure aumentare la riserva di fase del progetto continuo;
    \item il passaggio finale da \(C(s)\) a \(R(z)\) si fa con Tustin, in questo caso preferibilmente con pre-warping.
\end{enumerate}